\chapter{Beyond CS 136}

\textbf{Readings:} CP:AMA 17.7 (function pointers), 2.1 and 15.4 (compiling and linking),
26.2 (\code{qsort}).

\begin{aside}
Much of this unit is a preview of later courses (CS 246, CS 241, CS 240) rather than core exam
material; check with your instructor which parts are tested. Generic \code{void *} ADTs and
function pointers were the parts we actually used.
\end{aside}

\section{Generic Data with \code{void *}}

All our ADTs so far store \code{int}s. To store strings we would copy the whole module and change
the types. Racket does not have this problem because it is dynamically typed. In C, the
\emph{generic pointer} \code{void *} can point to data of any type:
\begin{itemize}
  \item any pointer can be assigned to a \code{void *} and back without a cast;
  \item a \code{void *} cannot be dereferenced (the type, and hence the size, is unknown) and
    pointer arithmetic on it is not allowed;
  \item the compiler can no longer check types, so the client must be careful.
\end{itemize}

\subsection{A generic stack}
\begin{lstlisting}[style=c]
// gstack.h: a generic stack; it stores pointers provided by the client.
//   The client owns the items: the stack never frees them.
struct stack;
struct stack *stack_create(void);
void stack_destroy(struct stack *s);   // requires: s is empty
bool stack_is_empty(const struct stack *s);
void stack_push(void *item, struct stack *s);
void *stack_top(const struct stack *s);
void *stack_pop(struct stack *s);
\end{lstlisting}
\begin{lstlisting}[style=c]
// gstack.c (dynamic array of void *)
struct stack {
  int len;
  int maxlen;
  void **data;              // array of generic pointers
};

void stack_push(void *item, struct stack *s) {
  assert(s);
  if (s->len == s->maxlen) {
    s->maxlen *= 2;
    s->data = realloc(s->data, s->maxlen * sizeof(void *));
  }
  s->data[s->len] = item;
  ++s->len;
}
\end{lstlisting}
\begin{lstlisting}[style=c]
// client: a stack of heap-allocated ints
struct stack *s = stack_create();
int *p = malloc(sizeof(int));
*p = 42;
stack_push(p, s);
int *q = stack_pop(s);        // void * converted back to int *
printf("%d\n", *q);           // 42
free(q);                      // the client frees its own items
stack_destroy(s);
\end{lstlisting}

\begin{important}[Ownership of generic items]
With \code{void *} ADTs the interface must say who frees the items. Common choices: the client
frees them (as above), or the client passes a ``destroy'' function pointer to
\code{stack\_destroy} so the ADT can free any remaining items.
\end{important}

\section{Function Pointers for Generic Algorithms}

A generic data structure that needs to \emph{compare} items (a sorted list, a BST dictionary)
cannot use \code{<} on \code{void *}. The client provides a comparison function instead, using
the same convention as \code{strcmp}: negative, zero or positive.

\subsection{\code{qsort}}

The C library's generic sort (in \code{<stdlib.h>}):
\begin{lstlisting}[style=c]
void qsort(void *base, size_t len, size_t size,
           int (*compare)(const void *, const void *));
\end{lstlisting}
\begin{lstlisting}[style=c]
// compare_ints(a, b) compares the ints a and b point at
int compare_ints(const void *a, const void *b) {
  const int *ia = a;
  const int *ib = b;
  if (*ia < *ib) return -1;
  if (*ia > *ib) return 1;
  return 0;
  // (return *ia - *ib; is shorter but can overflow)
}

int a[5] = {5, 3, 9, 1, 7};
qsort(a, 5, sizeof(int), compare_ints);   // a is {1, 3, 5, 7, 9}
\end{lstlisting}
Sorting an array of strings (\code{char *}) works the same way: each element is a
\code{char *}, so the comparator receives \code{char **}s:
\begin{lstlisting}[style=c]
int compare_strs(const void *a, const void *b) {
  const char * const *sa = a;
  const char * const *sb = b;
  return strcmp(*sa, *sb);
}
\end{lstlisting}

\subsection{A generic BST dictionary}
\begin{lstlisting}[style=c]
struct dictionary {
  struct bstnode *root;
  int (*key_compare)(const void *, const void *);
};

struct dictionary *dict_create(int (*cmp)(const void *, const void *)) {
  struct dictionary *d = malloc(sizeof(struct dictionary));
  d->root = NULL;
  d->key_compare = cmp;
  return d;
}

void *dict_lookup(const void *key, const struct dictionary *d) {
  const struct bstnode *node = d->root;
  while (node) {
    int result = d->key_compare(key, node->key);
    if (result == 0) return node->val;
    node = (result < 0) ? node->left : node->right;
  }
  return NULL;
}
\end{lstlisting}
This is a limited form of what Racket gets for free with higher-order functions, and what C++
achieves with templates (CS 246).

\section{\code{typedef}}

\code{typedef} creates an alias for a type. CS 136 style mostly avoided it (so that pointers and
structures stay visible), but it is common in real code, especially for function pointer types:
\begin{lstlisting}[style=c]
typedef int (*CompareFn)(const void *, const void *);
typedef struct dictionary *Dictionary;   // hides the pointer: use carefully

Dictionary dict_create(CompareFn cmp);
\end{lstlisting}

\section{The Compilation Pipeline}

Seashell hides how a C program becomes an executable. Running \code{gcc} on several files does
four steps:
\begin{center}
\begin{tikzpicture}[box/.style={draw,rounded corners,align=center,minimum height=0.9cm,
  text width=2.2cm,font=\small},
  arr/.style={-{Stealth[length=2mm]},thick}]
  \node[box] (src) {\code{.c} / \code{.h}\\ source};
  \node[box,right=0.6cm of src] (pre) {preprocessor};
  \node[box,right=0.6cm of pre] (cc) {compiler\\ (+ assembler)};
  \node[box,right=0.6cm of cc] (ld) {linker};
  \node[box,below=0.8cm of ld] (exe) {executable};
  \draw[arr] (src) -- (pre);
  \draw[arr] (pre) -- (cc);
  \draw[arr] (cc) -- node[above,font=\scriptsize]{\code{.o}} (ld);
  \draw[arr] (ld) -- (exe);
\end{tikzpicture}
\end{center}
\begin{enumerate}
  \item \textbf{Preprocessing:} handles directives starting with \code{\#}: \code{\#include}
    pastes in header files, \code{\#define} does textual macro substitution, \code{\#ifdef}
    includes code conditionally. The result is a single ``translation unit'' per \code{.c} file.
  \item \textbf{Compiling:} each translation unit is checked (types, syntax) and translated to
    assembly language and then (by the \emph{assembler}) to machine code in an \emph{object file}
    (\code{.o}). References to functions in other modules are left unresolved. This is why a
    \code{.c} file only needs the \emph{declarations} from headers.
  \item \textbf{Linking:} the linker combines all object files and libraries (e.g.\ the C
    standard library for \code{printf}), resolves every reference to its definition, and
    produces the executable. ``Undefined reference'' errors and duplicate-definition errors come
    from this step.
\end{enumerate}
\begin{lstlisting}[style=bash]
gcc -std=c99 -Wall -c stack.c        # produces stack.o
gcc -std=c99 -Wall -c main.c         # produces main.o
gcc stack.o main.o -o prog           # link
./prog < test1.in
\end{lstlisting}
Separate compilation means that when one module changes, only that module needs to be
recompiled before relinking. Tools like \code{make} automate this (CS 246).

\subsection{Include guards}

If a header is included twice in one translation unit (e.g.\ directly and through another
header), structure definitions would be duplicated. Headers are therefore protected:
\begin{lstlisting}[style=c]
#ifndef STACK_H
#define STACK_H
// ... declarations ...
#endif
\end{lstlisting}

\subsection{Command-line arguments}

Outside Seashell, \code{main} can receive arguments typed on the command line:
\begin{lstlisting}[style=c]
int main(int argc, char *argv[]) {
  // argc: number of arguments (including the program name)
  // argv: array of strings; argv[0] is the program name
  for (int i = 1; i < argc; ++i) {
    printf("argument %d: %s\n", i, argv[i]);
  }
}
\end{lstlisting}

\section{C vs.\ Racket vs.\ Other Languages}

\begin{center}
\begin{tabular}{p{3.2cm}p{3.6cm}p{3.6cm}p{3.6cm}}
  \toprule
  & Racket & C & C++ / Java / Python \\
  \midrule
  Paradigm & functional (mostly) & imperative & imperative + object-oriented \\
  Typing & dynamic & static (weak: casts, \code{void *}) & static (C++, Java), dynamic (Python) \\
  Memory & garbage collected & manual \code{malloc}/\code{free} & RAII/destructors (C++), GC (Java, Python) \\
  Execution & interpreted / JIT & compiled to machine code & compiled (C++), bytecode VM (Java), interpreted (Python) \\
  Generic code & automatic & \code{void *}, function pointers & templates / generics / duck typing \\
  Safety & errors reported & undefined behaviour & exceptions, bounds checks (Java, Python) \\
  \bottomrule
\end{tabular}
\end{center}

\begin{itemize}
  \item \textbf{Garbage collection:} Racket, Java and Python automatically free memory that is
    no longer reachable, so leaks and dangling pointers (mostly) cannot happen, at some cost in
    run-time overhead and control. C gives full control and full responsibility.
  \item \textbf{Object-oriented programming} (CS 246) bundles data with the functions that operate
    on it (\emph{classes}), with access control (\code{private}) enforcing information hiding ---
    a language-level version of our opaque-structure ADTs. Constructors and destructors replace
    \code{\_create} and \code{\_destroy}.
  \item \textbf{Undefined behaviour} is C's trade-off for speed: no bounds checks, no overflow
    checks. Many security vulnerabilities (buffer overflows) come from it.
  \item Despite its age, C remains essential for operating systems, embedded systems, and
    performance-critical libraries, and its memory model underlies how every other language is
    implemented.
\end{itemize}

\section{Where the Ideas Go Next}

\begin{itemize}
  \item \textbf{CS 246} (Object-Oriented Software Development): C++, classes, RAII, templates,
    design patterns, \code{make}.
  \item \textbf{CS 241} (Foundations of Sequential Programs): how the compiler, assembler and
    linker actually work; MIPS assembly; how the stack and heap are managed.
  \item \textbf{CS 240} (Data Structures and Data Management): balanced trees, hash tables,
    priority queues, sorting lower bounds, and more rigorous analysis.
  \item \textbf{CS 251} (Computer Organization and Design): how memory, bits and machine code are
    realised in hardware.
  \item \textbf{CS 350} (Operating Systems): processes, virtual memory --- the segmentation faults
    we saw come from here.
\end{itemize}

\section{Summary}
\begin{itemize}
  \item \code{void *} allows generic ADTs; the interface must define ownership of items.
  \item Function pointers (comparators like \code{qsort}'s) make algorithms generic.
  \item Building a program: preprocess $\to$ compile/assemble to object files $\to$ link.
    Headers hold declarations; include guards prevent double inclusion.
  \item C trades safety for control and speed; other languages add garbage collection,
    objects and generics on top of the same memory model.
\end{itemize}
